\documentclass[
    language=english,
    doctype=article,
    institution=none,
]{../../omnilatex}

\RequirePackage{config/document-settings}
\addbibresource{../../bib/bibliography.bib}

\RequirePackage{tikz}
\RequirePackage{pgfplots}
\pgfplotsset{compat=1.18}
\RequirePackage{booktabs}
\RequirePackage{longtable}
\RequirePackage{tabularx}
\RequirePackage[ruled,vlined,linesnumbered]{algorithm2e}
\RequirePackage{minted}
\RequirePackage{cleveref}
\RequirePackage{subcaption}
\RequirePackage{float}
\RequirePackage{multirow}
\RequirePackage{xcolor}
\RequirePackage[most]{tcolorbox}
\RequirePackage{cancel}
\RequirePackage{siunitx}
\RequirePackage{chemformula}
\RequirePackage{physics}
\RequirePackage{tikz-cd}
\RequirePackage{makecell}

% ─────────────────────────────────────────────────────
% Theorem Environments
% ─────────────────────────────────────────────────────
\theoremstyle{definition}
\newtheorem{theorem}{Theorem}[section]
\newtheorem{lemma}[theorem]{Lemma}
\newtheorem{proposition}[theorem]{Proposition}
\newtheorem{corollary}[theorem]{Corollary}
\theoremstyle{definition}
\newtheorem{definition}{Definition}[section]
\newtheorem{example}{Example}[section]
\theoremstyle{remark}
\newtheorem{remark}{Remark}[section]

% ─────────────────────────────────────────────────────
% Custom Commands
% ─────────────────────────────────────────────────────
\newcommand{\R}{\mathbb{R}}
\newcommand{\N}{\mathbb{N}}
\newcommand{\Z}{\mathbb{Z}}
\newcommand{\C}{\mathbb{C}}
\DeclareMathOperator*{\argmin}{arg\,min}
\DeclareMathOperator*{\argmax}{arg\,max}
\DeclareMathOperator{\Tr}{Tr}
\DeclareMathOperator{\diag}{diag}
\DeclareMathOperator{\sgn}{sgn}
\DeclareMathOperator{\Softmax}{Softmax}
\DeclareMathOperator{\Concat}{Concat}
\newcommand{\mat}[1]{\mathbf{#1}}
\newcommand{\tensor}[2]{#1^{#2}}
\newcommand{\hyp}{\mathcal{H}}
\newcommand{\dataset}{\mathcal{D}}
\newcommand{\graph}{\mathcal{G}}
\DeclareMathOperator{\Var}{Var}
\DeclareMathOperator{\Cov}{Cov}
\newcommand{\inner}[2]{\langle #1, #2 \rangle}

\begin{document}

\maketitle

% ─────────────────────────────────────────────────────
% Structured Abstract
% ─────────────────────────────────────────────────────
\begin{abstract}
\noindent\textbf{Background.}
Knowledge graph reasoning requires models that can traverse multiple hops
across heterogeneous relation types, yet existing graph neural networks
suffer from over-smoothing and relational information loss beyond three
hops~\cite{goossensLaTeXCompanion1993}.

\noindent\textbf{Methods.}
We introduce HeteroGNN, a novel architecture that combines relation-aware
attention with adaptive neighbourhood sampling.  The model employs a
hierarchical message-passing scheme with residual connections at each hop
to preserve long-range signal~\cite{knuthTeXbook1986}.

\noindent\textbf{Results.}
On four benchmark knowledge graphs, HeteroGNN achieves up to
\(\SI{8.3}{\percent}\) improvement in Mean Reciprocal Rank over the
strongest baseline while reducing training time by
\(\SI{35}{\percent}\)~\cite{mckinneyPythonDataAnalysis2018}.

\noindent\textbf{Conclusions.}
Our results demonstrate that relation-aware attention combined with
adaptive sampling provides a principled solution to multi-hop reasoning
over heterogeneous graphs, opening new directions for scalable knowledge
graph completion.
\end{abstract}

\tableofcontents

% =====================================================
\section{Introduction}
\label{sec:intro}
% =====================================================

Knowledge graphs (KGs) encode structured real-world knowledge as
triples~\(\langle \text{head}, \text{relation}, \text{tail}\rangle\)
and power applications from question answering to drug
discovery~\cite{einstein1905}.  Multi-hop reasoning---inferring facts
that require traversing multiple edges---remains a fundamental challenge.

\textcite{knuthTeXbook1986} identified three core difficulties: (i) the
combinatorial explosion of reachable nodes, (ii) the heterogeneity of
relation types, and (iii) the degradation of learned representations at
greater depths.  Despite significant progress, no existing method
simultaneously addresses all three.

\subsection{Research Gap}

Existing graph neural networks (GNNs) for KG reasoning fall into two
categories.  Path-based methods~\cite{mollenhauerHandbuchDieselmotoren2007}
enumerate explicit paths but scale exponentially with hop count.
Message-passing methods~\cite{ramalhoFluentPython2015} aggregate
neighbourhood information but suffer from over-smoothing, where node
representations converge to indistinguishable vectors after several
layers~\cite{baehrThermodynamikGrundlagenUnd2016}.

\subsection{Contributions}

This paper makes the following contributions:
\begin{enumerate}
    \item We propose \emph{HeteroGNN}, a relation-aware graph neural
          network with adaptive neighbourhood sampling
          (\cref{sec:methodology}).
    \item We prove formal bounds on the information preservation across
          multiple message-passing hops (\cref{thm:info-preserve}).
    \item We conduct extensive experiments on four benchmarks, achieving
          state-of-the-art results (\cref{sec:experiments}).
\end{enumerate}

% =====================================================
\section{Related Work}
\label{sec:related}
% =====================================================

\subsection{Graph Neural Networks for Reasoning}

Early approaches applied spectral graph convolutions to
KGs~\cite{diracPrinciplesQuantumMechanics1981}.  More recent work
focuses on spatial message-passing, where each node aggregates
information from its neighbours.

\Cref{tab:related} compares the key properties of existing approaches.

\begin{table}[htbp]
    \centering
    \caption{Comparison of graph neural network architectures for knowledge
    graph reasoning.}
    \label{tab:related}
    \begin{tabular}{@{}llccccc@{}}
        \toprule
        \textbf{Method}
            & \textbf{Year}
            & \multicolumn{1}{c}{\textbf{Hetero.}}
            & \multicolumn{1}{c}{\textbf{Multi-hop}}
            & \multicolumn{1}{c}{\textbf{Attn.}}
            & \multicolumn{1}{c}{\textbf{Adapt.\ Sampling}}
            & \multicolumn{1}{c}{\textbf{Scalable}} \\
        \midrule
        \multirow{2}{*}{R-GCN}
            & 2018 & \checkmark & & & & \checkmark \\
            & & \multicolumn{5}{l}{\footnotesize Per-relation weight matrices} \\
        \midrule
        \multirow{2}{*}{CompGCN}
            & 2020 & \checkmark & & \checkmark & & \checkmark \\
            & & \multicolumn{5}{l}{\footnotesize Compositional embeddings} \\
        \midrule
        \multirow{2}{*}{HAN}
            & 2019 & \checkmark & & \checkmark & & \\
            & & \multicolumn{5}{l}{\footnotesize Hierarchical attention} \\
        \midrule
        \multirow{2}{*}{NBFNet}
            & 2021 & & \checkmark & & & \\
            & & \multicolumn{5}{l}{\footnotesize Bellman--Ford message passing} \\
        \midrule
        \multirow{2}{*}{\textbf{Ours}}
            & 2026 & \checkmark & \checkmark & \checkmark & \checkmark & \checkmark \\
            & & \multicolumn{5}{l}{\footnotesize Relation-aware attention + adaptive sampling} \\
        \bottomrule
    \end{tabular}
\end{table}

% =====================================================
\section{Methodology}
\label{sec:methodology}
% =====================================================

\subsection{Problem Definition}

\begin{definition}[Heterogeneous Knowledge Graph]
    \label{def:hkg}
    A heterogeneous knowledge graph is a tuple
    \(\graph = (\mathcal{V}, \mathcal{E}, \mathcal{R}, \phi)\)
    where \(\mathcal{V}\) is a set of nodes,
    \(\mathcal{E} \subseteq \mathcal{V} \times \mathcal{R} \times \mathcal{V}\)
    is a set of edges, \(\mathcal{R}\) is a finite set of relation types,
    and \(\phi \colon \mathcal{V} \to \mathcal{T}\) maps nodes to types.
\end{definition}

The goal of multi-hop link prediction is: given a query triple
\((h, r, ?)\), predict the tail entity \(t\) by reasoning over paths of
length \(k\) in \(\graph\).

\subsection{Algorithm}

\Cref{alg:heternn} describes the forward pass of HeteroGNN.

\begin{algorithm}[htbp]
    \caption{HeteroGNN Forward Pass}
    \label{alg:heternn}
    \SetAlgoLined
    \SetKwInOut{Input}{Input}
    \SetKwInOut{Output}{Output}

    \Input{Graph \(\graph\), node features \(\mat{X} \in \R^{|\mathcal{V}|\times d}\),
           number of layers \(L\), sample size \(S\)}
    \Output{Updated node representations \(\mat{H}^{(L)}\)}

    \(\mat{H}^{(0)} \gets \mat{X}\) \tcp*{initialise}
    \For{$\ell \gets 1$ \KwTo $L$}{
        \ForEach{relation $r \in \mathcal{R}$}{
            \(\mathcal{N}_r(v) \gets\) \textsc{AdaptiveSample}
                \((\graph, v, r, S)\) \tcp*{\(\mathcal{O}(|\mathcal{E}_r|)\)}
            \(\mat{m}_r^{(\ell)} \gets
                \sum_{u \in \mathcal{N}_r(v)}
                \alpha_{vu}^{r}\,
                \mat{W}_r^{(\ell)}\,\mat{h}_u^{(\ell-1)}\)
                \tcp*{relation-aware message, \(\mathcal{O}(S\,d^2)\)}
        }
        \(\mat{h}_v^{(\ell)} \gets
            \sigma\!\left(
                \mat{W}_{\text{res}}^{(\ell)}\,\mat{h}_v^{(\ell-1)}
                + \sum_{r \in \mathcal{R}}
                \beta_r\,\mat{m}_r^{(\ell)}
            \right)\)
            \tcp*{residual update, \(\mathcal{O}(|\mathcal{R}|\,d)\)}
    }
    \Return \(\mat{H}^{(L)}\)\;
\end{algorithm}

\subsection{Complexity Analysis}

\begin{proposition}[Time Complexity]
    \label{prop:complexity}
    The total time complexity of HeteroGNN per forward pass is
    \(\mathcal{O}\bigl(L \cdot |\mathcal{R}| \cdot (S\,d^2 + |\mathcal{V}|\,d)\bigr)\),
    where \(L\) is the number of layers, \(|\mathcal{R}|\) the number of
    relation types, \(S\) the sample size, and \(d\) the hidden dimension.
\end{proposition}

\subsection{Mathematical Proof}

\begin{theorem}[Information Preservation]
    \label{thm:info-preserve}
    Let \(\mat{H}^{(\ell)}\) be the node representations at layer \(\ell\).
    If the residual connection satisfies
    \(\norm{\mat{W}_{\text{res}}^{(\ell)} - \mat{I}} \le \epsilon < 1\),
    then the mutual information between input features and layer-\(L\)
    representations satisfies
    \(\mathcal{I}\bigl(\mat{X};\,\mat{H}^{(L)}\bigr)
        \ge (1-\epsilon)^L \cdot \mathcal{I}\bigl(\mat{X};\,\mat{H}^{(0)}\bigr)\).
\end{theorem}

\begin{proof}
    By induction on \(L\).  For \(L=0\), the bound holds trivially.
    Assume it holds for \(L-1\).  At layer \(L\), the residual connection
    gives
    \[
        \mat{h}_v^{(L)}
            = \mat{W}_{\text{res}}^{(L)}\,\mat{h}_v^{(L-1)}
              + \mat{f}^{(L)}\!\bigl(\mathcal{N}(v)\bigr),
    \]
    where \(\mat{f}^{(L)}\) is the aggregated message.  By the data
    processing inequality applied to the linear map
    \(\mat{W}_{\text{res}}^{(L)}\),
    \begin{align}
        \mathcal{I}\bigl(\mat{X};\,\mat{H}^{(L)}\bigr)
            &\ge \mathcal{I}\bigl(\mat{X};\,
               \mat{W}_{\text{res}}^{(L)}\,\mat{H}^{(L-1)}\bigr) \\
            &\ge (1-\epsilon)\,\mathcal{I}\bigl(\mat{X};\,\mat{H}^{(L-1)}\bigr) \\
            &\ge (1-\epsilon)^L\,\mathcal{I}\bigl(\mat{X};\,\mat{H}^{(0)}\bigr),
    \end{align}
    where the last step applies the inductive hypothesis.
\end{proof}

\begin{corollary}
    \label{cor:residual-depth}
    With \(\epsilon = 0.1\), at least \(\SI{65}{\percent}\) of the input
    information is preserved after \(L=5\) layers.
\end{corollary}

\subsection{System Architecture}

\Cref{fig:arch} shows the overall architecture of HeteroGNN.

\begin{figure}[htbp]
    \centering
    \begin{tikzpicture}[
        node distance=1.4cm, auto, >=stealth',
        block/.style={rectangle, draw, fill=blue!15, text width=2.5cm,
                      text centered, rounded corners, minimum height=0.8cm},
        io/.style={trapezium, trapezium left angle=70,
                   trapezium right angle=110, draw, fill=green!15,
                   text width=2cm, text centered, minimum height=0.8cm},
        arr/.style={->, thick},
    ]
        \node[io]    (input)                {Input KG};
        \node[block, right of=input, xshift=1.8cm]  (embed)  {Relation\\Embeddings};
        \node[block, right of=embed, xshift=1.8cm]  (sample) {Adaptive\\Sampler};
        \node[block, right of=sample, xshift=1.8cm] (agg)    {Relation-Aware\\Aggregation};
        \node[block, below of=agg, yshift=-0.8cm]   (res)    {Residual\\Connection};
        \node[io,    below of=res, yshift=-0.8cm]   (out)    {Predictions};

        \draw[arr] (input)  -- (embed);
        \draw[arr] (embed)  -- (sample);
        \draw[arr] (sample) -- (agg);
        \draw[arr] (agg)    -- (res);
        \draw[arr] (res)    -- (out);
        \draw[arr, dashed] (res.west) -- ++(-2,0) |- node[near start, left] {skip} (sample.south);
    \end{tikzpicture}
    \caption{Architecture of the HeteroGNN model showing the main
    computational stages.}
    \label{fig:arch}
\end{figure}

\subsection{Data Flow}

\Cref{fig:dataflow} illustrates how data transforms through the pipeline.

\begin{figure}[htbp]
    \centering
    \begin{tikzpicture}[>=stealth', node distance=2.2cm, auto,
        proc/.style={rectangle, draw, fill=orange!15, text width=2.2cm,
                     text centered, minimum height=0.7cm, rounded corners},
    ]
        \node[proc] (raw)   {Raw Triples};
        \node[proc, right of=raw, xshift=1.5cm]   (enc)   {Tokenise};
        \node[proc, right of=enc, xshift=1.5cm]   (feat)  {Feature Map};
        \node[proc, right of=feat, xshift=1.5cm]  (gnn)   {HeteroGNN};
        \node[proc, below of=gnn, yshift=-0.8cm]  (dec)   {Decode};

        \draw[->, thick, cBlue]   (raw)  -- node {triples} (enc);
        \draw[->, thick, cGreen]  (enc)  -- node {\(\mat{X}\)} (feat);
        \draw[->, thick, cRed]    (feat) -- node {context} (gnn);
        \draw[->, thick, violet]  (gnn)  -- node {\(\mat{H}\)} (dec);
    \end{tikzpicture}
    \caption{Data flow from raw knowledge graph triples to predictions.}
    \label{fig:dataflow}
\end{figure}

\subsection{Neural Network Diagram}

\Cref{fig:nn} visualises the layer structure of HeteroGNN.

\begin{figure}[htbp]
    \centering
    \begin{tikzpicture}[
        neuron/.style={circle, draw=black!80, thick, minimum size=20pt,
                       inner sep=0pt, font=\footnotesize},
        in/.style={neuron, fill=cBlue!20},
        hid/.style={neuron, fill=cGreen!20},
        out/.style={neuron, fill=cRed!20},
        lbl/.style={text width=3em, align=center, font=\footnotesize},
    ]
        \foreach \i in {1,...,4}
            \node[in]  (I-\i) at (0, -\i*1.1) {$e_{\i}$};
        \foreach \j in {1,...,5}
            \node[hid] (H-\j) at (2.5, -\j*1.1+0.55) {$h_{\j}$};
        \foreach \k in {1,...,3}
            \node[out] (O-\k) at (5, -\k*1.1-0.55) {$o_{\k}$};

        \foreach \i in {1,...,4}
            \foreach \j in {1,...,5}
                \draw[->, thin, gray!50] (I-\i) -- (H-\j);
        \foreach \j in {1,...,5}
            \foreach \k in {1,...,3}
                \draw[->, thin, gray!50] (H-\j) -- (O-\k);

        \node[lbl, above=10pt] at (0,0)   {Input};
        \node[lbl, above=10pt] at (2.5,0) {Hidden};
        \node[lbl, above=10pt] at (5,0)   {Output};
    \end{tikzpicture}
    \caption{Layer structure of HeteroGNN with 4 input, 5 hidden, and
    3 output units.}
    \label{fig:nn}
\end{figure}

% =====================================================
\section{Experiments}
\label{sec:experiments}
% =====================================================

\subsection{Experimental Setup}

\Cref{tab:setup} summarises the datasets used.

\begin{table}[htbp]
    \centering
    \caption{Dataset statistics for the four knowledge graph benchmarks.}
    \label{tab:setup}
    \begin{tabular}{@{}lrrrr@{}}
        \toprule
        \textbf{Dataset}
            & \textbf{\#Nodes}
            & \textbf{\#Edges}
            & \textbf{\#Relations}
            & \textbf{\#Train / Val / Test} \\
        \midrule
        WN18RR   & 40\,943   & 93\,003   & 11 & 86\,835 / 3\,034 / 3\,134 \\
        FB15K-237 & 14\,541  & 310\,116  & 237 & 272\,115 / 17\,535 / 20\,466 \\
        YAGO3-10 & 123\,182 & 1\,079\,040 & 37 & 936\,090 / 22\,826 / 20\,130 \\
        NELL-995 & 75\,492  & 200\,423  & 200 & 139\,700 / 14\,968 / 14\,968 \\
        \bottomrule
    \end{tabular}
\end{table}

\subsection{Main Results}

\Cref{fig:main-results} shows the primary performance comparison across
all datasets.

\begin{figure}[htbp]
    \centering
    \begin{tikzpicture}
    \begin{axis}[
        width=0.85\textwidth,
        height=6.5cm,
        ybar,
        bar width=12pt,
        ylabel={MRR},
        symbolic x coords={WN18RR, FB15K-237, YAGO3-10, NELL-995},
        xtick=data,
        ymin=0.1, ymax=0.6,
        legend style={at={(0.5,-0.2)}, anchor=north, legend columns=4,
                      font=\footnotesize},
        grid=major,
        grid style={dashed, gray!30},
        nodes near coords,
        every node near coord/.append style={font=\tiny, rotate=90,
                                              anchor=west},
    ]
        \addplot[fill=cBlue!60] coordinates
            {(WN18RR,0.42) (FB15K-237,0.30) (YAGO3-10,0.35) (NELL-995,0.28)};
        \addplot[fill=cRed!60] coordinates
            {(WN18RR,0.47) (FB15K-237,0.34) (YAGO3-10,0.39) (NELL-995,0.31)};
        \addplot[fill=cGreen!60] coordinates
            {(WN18RR,0.49) (FB15K-237,0.36) (YAGO3-10,0.41) (NELL-995,0.33)};
        \addplot[fill=violet!60] coordinates
            {(WN18RR,0.53) (FB15K-237,0.41) (YAGO3-10,0.46) (NELL-995,0.38)};
        \addlegendentry{R-GCN}
        \addlegendentry{CompGCN}
        \addlegendentry{NBFNet}
        \addlegendentry{\textbf{HeteroGNN}}
    \end{axis}
    \end{tikzpicture}
    \caption{Mean Reciprocal Rank (MRR) on four knowledge graph benchmarks.
    Higher is better.}
    \label{fig:main-results}
\end{figure}

\subsection{Ablation Study}

\Cref{fig:ablation} isolates the contribution of each component.

\begin{figure}[htbp]
    \centering
    \begin{tikzpicture}
    \begin{groupplot}[
        group style={
            group size=2 by 1,
            horizontal sep=2.5cm,
        },
        width=0.42\textwidth,
        height=6cm,
        ylabel={MRR},
        symbolic x coords={Full, No Attn, No Res, No Sample},
        xtick=data,
        x tick label style={rotate=30, anchor=east, font=\footnotesize},
        ybar=4pt,
        bar width=10pt,
        legend style={at={(0.5,-0.3)}, anchor=north, legend columns=-1,
                      font=\footnotesize},
        ymin=0.3, ymax=0.6,
        grid=major,
        grid style={dashed, gray!40},
    ]
        \nextgroupplot[title={\textbf{WN18RR}}]
            \addplot[fill=cBlue!70] coordinates
                {(Full,0.53) (No Attn,0.47) (No Res,0.44) (No Sample,0.51)};
            \addplot[fill=cRed!70] coordinates
                {(Full,0.51) (No Attn,0.45) (No Res,0.42) (No Sample,0.49)};
            \addlegendentry{3 hops}
            \addlegendentry{5 hops}

        \nextgroupplot[title={\textbf{FB15K-237}}]
            \addplot[fill=cBlue!70] coordinates
                {(Full,0.41) (No Attn,0.36) (No Res,0.33) (No Sample,0.39)};
            \addplot[fill=cRed!70] coordinates
                {(Full,0.39) (No Attn,0.34) (No Res,0.31) (No Sample,0.37)};
            \addlegendentry{3 hops}
            \addlegendentry{5 hops}
    \end{groupplot}
    \end{tikzpicture}
    \caption{Ablation study: removing attention, residual connections, or
    adaptive sampling degrades performance at both 3 and 5 hops.}
    \label{fig:ablation}
\end{figure}

\subsection{Scalability Analysis}

\Cref{fig:scalability} examines the trade-off between model size and
training time.

\begin{figure}[htbp]
    \centering
    \begin{tikzpicture}
    \begin{axis}[
        width=0.85\textwidth,
        height=6cm,
        xlabel={Number of Parameters [\si{\mega}]},
        ylabel={MRR},
        axis y line*=left,
        grid=major,
        thick,
        legend pos=north west,
    ]
        \addplot[cBlue, thick, mark=*, mark size=2.5pt] coordinates
            {(2.1, 0.42) (4.3, 0.47) (8.7, 0.51) (17.4, 0.53) (34.8, 0.54)};
        \addlegendentry{MRR}
    \end{axis}
    \begin{axis}[
        width=0.85\textwidth,
        height=6cm,
        ylabel={Time [\si{\hour}]},
        axis y line*=right,
        axis x line=none,
        thick,
        legend pos=north east,
    ]
        \addplot[cRed, thick, dashed, mark=square*, mark size=2.5pt] coordinates
            {(2.1, 1.2) (4.3, 2.8) (8.7, 5.1) (17.4, 10.3) (34.8, 22.5)};
        \addlegendentry{Training time}
    \end{axis}
    \end{tikzpicture}
    \caption{Scalability analysis: MRR improves with model size (left axis)
    at the cost of training time (right axis).}
    \label{fig:scalability}
\end{figure}

\subsection{Failure Analysis}

\Cref{fig:failure} categorises the types of errors made by HeteroGNN.

\begin{figure}[htbp]
    \centering
    \begin{tikzpicture}[font=\footnotesize]
        \def\radius{2.5}
        \draw[fill=cBlue!70,  draw=white, thick] (0,0) -- (0:\radius)
            arc(0:144:\radius) -- cycle;
        \draw[fill=cRed!70,   draw=white, thick] (0,0) -- (144:\radius)
            arc(144:234:\radius) -- cycle;
        \draw[fill=cGreen!70, draw=white, thick] (0,0) -- (234:\radius)
            arc(234:306:\radius) -- cycle;
        \draw[fill=orange!70, draw=white, thick] (0,0) -- (306:\radius)
            arc(306:360:\radius) -- cycle;

        \node[font=\bfseries] at (72:{0.65*\radius})  {40\%};
        \node[font=\bfseries] at (189:{0.65*\radius}) {25\%};
        \node[font=\bfseries] at (270:{0.65*\radius}) {20\%};
        \node[font=\bfseries] at (333:{0.65*\radius}) {15\%};

        \fill[cBlue!70]  (\radius+0.6,  0.0) rectangle ++(0.3,0.3);
        \node[anchor=west] at (\radius+1.0, 0.15) {Relation ambiguity};
        \fill[cRed!70]   (\radius+0.6, -0.5) rectangle ++(0.3,0.3);
        \node[anchor=west] at (\radius+1.0,-0.35) {Long-tail entities};
        \fill[cGreen!70] (\radius+0.6, -1.0) rectangle ++(0.3,0.3);
        \node[anchor=west] at (\radius+1.0,-0.85) {Noisy triples};
        \fill[orange!70] (\radius+0.6, -1.5) rectangle ++(0.3,0.3);
        \node[anchor=west] at (\radius+1.0,-1.35) {Missing links};
    \end{tikzpicture}
    \caption{Distribution of error types in the failure analysis.}
    \label{fig:failure}
\end{figure}

\subsection{Qualitative Results}

\Cref{fig:qualitative} shows example predictions on multi-hop queries.

\begin{figure}[htbp]
    \centering
    \subcaptionbox{2-hop query\label{fig:qual-2hop}}[0.45\textwidth]{
        \begin{tikzpicture}[>=stealth', node distance=1.2cm,
            ent/.style={circle, draw, fill=cBlue!15, minimum size=24pt,
                        inner sep=0pt, font=\small},
            rel/.style={font=\footnotesize, sloped, above},
        ]
            \node[ent] (h) {$h$};
            \node[ent, right of=h, xshift=1.5cm] (m) {$m$};
            \node[ent, right of=m, xshift=1.5cm] (t) {$t$};
            \draw[->, thick] (h) -- node[rel] {$r_1$} (m);
            \draw[->, thick] (m) -- node[rel] {$r_2$} (t);
        \end{tikzpicture}
    }
    \hfill
    \subcaptionbox{3-hop query\label{fig:qual-3hop}}[0.45\textwidth]{
        \begin{tikzpicture}[>=stealth', node distance=1.2cm,
            ent/.style={circle, draw, fill=cGreen!15, minimum size=24pt,
                        inner sep=0pt, font=\small},
            rel/.style={font=\footnotesize, sloped, above},
        ]
            \node[ent] (h) {$h$};
            \node[ent, right of=h, xshift=1.2cm] (m1) {$m_1$};
            \node[ent, right of=m1, xshift=1.2cm] (m2) {$m_2$};
            \node[ent, right of=m2, xshift=1.2cm] (t) {$t$};
            \draw[->, thick] (h)  -- node[rel] {$r_1$} (m1);
            \draw[->, thick] (m1) -- node[rel] {$r_2$} (m2);
            \draw[->, thick] (m2) -- node[rel] {$r_3$} (t);
        \end{tikzpicture}
    }
    \caption{Multi-hop query examples: (\subref{fig:qual-2hop}) a 2-hop
    reasoning path and (\subref{fig:qual-3hop}) a 3-hop path.}
    \label{fig:qualitative}
\end{figure}

% =====================================================
\section{Discussion}
\label{sec:discussion}
% =====================================================

\subsection{Key Findings}

Our experiments demonstrate that relation-aware attention consistently
outperforms relation-agnostic aggregation, confirming that heterogeneous
edge types carry complementary signal.  The residual connection is
critical for multi-hop reasoning: without it, performance drops sharply
at five hops (\cref{fig:ablation}).

\subsection{Limitations}

Several limitations merit attention.  First, the adaptive sampling
strategy assumes a uniform distribution of relation types, which may not
hold in highly imbalanced KGs.  Second, our theoretical bound
(\cref{thm:info-preserve}) is worst-case; tighter analysis using mutual
information estimators may yield sharper guarantees.  Third, we have not
explored dynamic graph settings where edges arrive over time.

\subsection{Future Work}

Promising directions include: (i) combining HeteroGNN with pre-trained
language models for joint reasoning over text and graph
structures~\cite{mollenhauerHandbuchDieselmotoren2007};
(ii) extending to temporal knowledge graphs where edge timestamps
encode evolution; and (iii) exploring federated training across
distributed KG partitions for privacy-preserving applications.

% =====================================================
\section{Conclusion}
\label{sec:conclusion}
% =====================================================

This paper introduced HeteroGNN, a relation-aware graph neural network
for multi-hop reasoning over heterogeneous knowledge graphs.  Our key
contributions are: (1) a relation-aware attention mechanism that preserves
heterogeneous structure, (2) a theoretical bound on information
preservation via residual connections (\cref{thm:info-preserve}), and
(3) extensive experiments demonstrating state-of-the-art performance on
four benchmarks (\cref{sec:experiments}).

% =====================================================
\section{Supplementary Material}
\label{sec:supplementary}
% =====================================================

\subsection{Extended Results}

\Cref{tab:extended} reports Hits@\(k\) metrics beyond MRR.

\begin{table}[htbp]
    \centering
    \caption{Extended results: Hits@1, Hits@3, and Hits@10 on all benchmarks.}
    \label{tab:extended}
    \begin{tabular}{@{}llcccc@{}}
        \toprule
        \textbf{Dataset} & \textbf{Method}
            & \textbf{H@1} & \textbf{MRR}
            & \textbf{H@3} & \textbf{H@10} \\
        \midrule
        WN18RR
            & R-GCN   & 0.34 & 0.42 & 0.47 & 0.56 \\
            & CompGCN & 0.38 & 0.47 & 0.52 & 0.61 \\
            & \textbf{Ours} & \textbf{0.45} & \textbf{0.53} & \textbf{0.58} & \textbf{0.67} \\
        \midrule
        FB15K-237
            & R-GCN   & 0.22 & 0.30 & 0.34 & 0.44 \\
            & CompGCN & 0.25 & 0.34 & 0.38 & 0.48 \\
            & \textbf{Ours} & \textbf{0.32} & \textbf{0.41} & \textbf{0.46} & \textbf{0.56} \\
        \midrule
        YAGO3-10
            & R-GCN   & 0.27 & 0.35 & 0.39 & 0.48 \\
            & CompGCN & 0.30 & 0.39 & 0.43 & 0.52 \\
            & \textbf{Ours} & \textbf{0.38} & \textbf{0.46} & \textbf{0.50} & \textbf{0.59} \\
        \midrule
        NELL-995
            & R-GCN   & 0.20 & 0.28 & 0.32 & 0.41 \\
            & CompGCN & 0.23 & 0.31 & 0.35 & 0.44 \\
            & \textbf{Ours} & \textbf{0.30} & \textbf{0.38} & \textbf{0.42} & \textbf{0.52} \\
        \bottomrule
    \end{tabular}
\end{table}

\subsection{Implementation}

\Cref{lst:python} shows the core aggregation function.

\begin{listing}[htbp]
    \centering
    \begin{minted}[
        linenos,
        stepnumber=5,
        fontsize=\small,
        frame=leftline,
        framerule=1pt,
        rulecolor=\color{g3},
        breaklines,
    ]{python}
import torch
import torch.nn as nn
from torch_geometric.nn import MessagePassing

class HeteroConvLayer(MessagePassing):
    def __init__(self, in_dim, out_dim, num_relations):
        super().__init__(aggr="add")
        self.W = nn.ModuleList([
            nn.Linear(in_dim, out_dim) for _ in range(num_relations)
        ])
        self.residual = nn.Linear(in_dim, out_dim)
        self.attn = nn.MultiheadAttention(out_dim, num_heads=4)

    def forward(self, x, edge_index_dict):
        msgs = []
        for r, (edge_index, _) in edge_index_dict.items():
            h = self.W[r](x)
            msg = self.propagate(edge_index, x=h)
            msgs.append(msg)
        agg = torch.stack(msgs).mean(dim=0)
        res = self.residual(x)
        return torch.relu(res + agg)
    \end{minted}
    \caption{Core HeteroGNN layer in PyTorch Geometric.}
    \label{lst:python}
\end{listing}

% ─────────────────────────────────────────────────────
% Cross-References
% ─────────────────────────────────────────────────────
\subsection{Cross-Reference Examples}

OmniLaTeX's \texttt{cleveref} integration provides intelligent
cross-references throughout:

\begin{itemize}
    \item \verb|\cref{thm:info-preserve}| \(\to\) \cref{thm:info-preserve}
    \item \verb|\cref{alg:heternn}| \(\to\) \cref{alg:heternn}
    \item \verb|\cref{fig:main-results}| \(\to\) \cref{fig:main-results}
    \item \verb|\cref{tab:extended}| \(\to\) \cref{tab:extended}
    \item \verb|\cref{def:hkg}| \(\to\) \cref{def:hkg}
    \item \verb|\cref{lst:python}| \(\to\) \cref{lst:python}
    \item \verb|\cref{fig:qual-2hop,fig:qual-3hop}|
          \(\to\) \cref{fig:qual-2hop,fig:qual-3hop}
\end{itemize}

% ─────────────────────────────────────────────────────
% Bibliography
% ─────────────────────────────────────────────────────
\printbibliography

\end{document}
